\section{Research Experience}

\entry{Medical University of Innsbruck}{Innsbruck, Austria}{PhD Researcher in adaptive optics}{Jan 2022 to Jun 2026}
\begin{adjustwidth}{103pt}{0pt}
\sepbullet \textbf{End to End System Design:} Led the hardware to software
implementation of multiphoton microscope, integrating laser sources, galvo scanners,
and light modulators. Achieved enhancing optical image quality.
\sepbullet \textbf{Systems Engineering:} instrument development following
V-model principles, managing requirements and tracking from concept to component
testing.
\sepbullet \textbf{Quality \& Testing:} Executed Design of Experiments
(DoE) and optical simulations in Python. Applied testing and documentation
protocols for optomechanical assemblies
frameworks.
\sepbullet \textbf{CAD \& Prototyping:} Designed custom optomechanical mounts and
alignment tools using SolidWorks and FreeCAD, ensuring precise mechanical
tolerances and seamless integration.
\end{adjustwidth}

\entry{Leibniz Institute of Photonic Technology}{Jena, Germany}{R\&D Research Assistant}{Mar 2020 to Dec 2021}
\begin{adjustwidth}{103pt}{0pt}
\sepbullet \textbf{Spectrometer Testing:} Testing and modifying a dispersive 1064
nm fiber probe Raman spectrometer for non-invasive human tissue characterization,
handling optomechanical layout and detector calibration.
\sepbullet \textbf{Data Analysis \& ML:} Developed a custom Python package
(SpectraMap) to apply supervised machine learning algorithms to 2D hyperspectral
Raman images, achieving high accuracy classification of tissue samples.
\sepbullet \textbf{Technical Documentation:} Authored peer reviewed publications detailing optical design
parameters and calibration procedures.
\end{adjustwidth}
